\editionPage[\pdfbookmark[0]{Démonstrations}{bm-preuves-debut-parent}]{preuves-debut}{Comprendre les preuves}{Une identité et une parité}{Les arguments ci-dessous éclairent des propriétés utilisées dans le livre.}
\begin{blocEdition}{preuves-debut}

\begin{demonstration}[Différence de deux carrés]
Par double distributivité, $(a-b)(a+b)=a^2+ab-ab-b^2=a^2-b^2$.
Les deux termes opposés s’annulent. L’égalité permet aussi bien de
développer le produit que de factoriser la différence.
\end{demonstration}
\begin{demonstration}[Le carré d’un entier impair est impair]
Un entier impair s’écrit $2n+1$ avec n entier. Son carré vaut
$(2n+1)^2=4n^2+4n+1=2(2n^2+2n)+1$. Le nombre entre parenthèses est entier :
on obtient bien un multiple de 2 augmenté de 1.
\end{demonstration}
\editionRetenir{Une écriture générale prouve ce que quelques essais ne font que suggérer.}

\end{blocEdition}
\editionPage{preuves-2}{Comprendre les preuves}{Pythagore par les aires}{Les arguments ci-dessous éclairent des propriétés utilisées dans le livre.}
\begin{blocEdition}{preuves-2}

Considérons quatre copies d’un triangle rectangle de côtés perpendiculaires
a et b et d’hypoténuse c, disposées dans un carré de côté $a+b$.
\begin{center}\begin{tikzpicture}[scale=.65,font=\small]
\draw[geoTrait](0,0)rectangle(5,5);
\fill[otGreenTint](2,0)--(5,2)--(3,5)--(0,3)--cycle;
\draw[geoTrait](2,0)--(5,2)--(3,5)--(0,3)--cycle;
\node[below]at(1,0){$a$};\node[below]at(3.5,0){$b$};
\node[right]at(5,1){$a$};\node[right]at(5,3.5){$b$};
\node at(2.5,2.5){aire $c^2$};
\end{tikzpicture}\end{center}
Chaque côté du quadrilatère central est une hypoténuse de longueur c.
Ses angles sont droits : les deux angles aigus voisins des triangles
ont une somme de 90 degrés, et complètent un angle plat avec l’angle central.
Le quadrilatère est donc un carré.\par\medskip
L’aire du grand carré vaut $(a+b)^2$. Elle est aussi la somme de l’aire
du carré central et de celles des quatre triangles :
\[
(a+b)^2=c^2+4\times\frac{ab}{2}=c^2+2ab.
\]
En développant puis en soustrayant $2ab$, on obtient
$\res{a^2+b^2=c^2}$.

\end{blocEdition}
\editionPage{preuves-3}{Comprendre les preuves}{Angles et triangles semblables}{Les arguments ci-dessous éclairent des propriétés utilisées dans le livre.}
\begin{blocEdition}{preuves-3}

\begin{demonstration}[La somme des angles d’un triangle]
Par le sommet C du triangle ABC, tracer la parallèle à AB.
Les angles alternes-internes y reproduisent les angles A et B,
de part et d’autre de l’angle C. Les trois forment un angle plat :
leur somme vaut $180^\circ$.
\end{demonstration}
\begin{center}\begin{tikzpicture}[scale=.75,font=\small]
\draw[geoTrait](0,0)node[below]{$A$}--(5,0)node[below]{$B$}--(2,3)node[above]{$C$}--cycle;
\draw[geoAide](-1,3)--(6,3);
\node at(1.1,2.7){$\alpha$};\node at(2,2.3){$\gamma$};\node at(3.1,2.7){$\beta$};
\end{tikzpicture}\end{center}
\begin{demonstration}[Deux angles égaux suffisent]
Dans deux triangles, si deux angles sont égaux deux à deux,
le troisième est chaque fois $180^\circ$ moins la même somme.
Il est donc égal aussi : les triangles sont semblables.
\end{demonstration}

\end{blocEdition}
\editionPage{preuves-4}{Comprendre les preuves}{Facteurs d’aire et de volume}{Les arguments ci-dessous éclairent des propriétés utilisées dans le livre.}
\begin{blocEdition}{preuves-4}

\begin{demonstration}[Agrandir un rectangle]
Un rectangle de côtés L et $\ell$ devient un rectangle de côtés
$kL$ et $k\ell$ lorsque $k>0$. Son aire devient
$(kL)(k\ell)=k^2L\ell$. Une figure découpée en triangles suit le même
principe, en multipliant base et hauteur par k. Pour une homothétie
de rapport négatif, la longueur est multipliée par $|k|$, et
$|k|^2=k^2$.
\end{demonstration}
\begin{demonstration}[Agrandir un pavé]
Le volume $L\ell h$ devient $(|k|L)(|k|\ell)(|k|h)=|k|^3L\ell h$.
Ce calcul explique pourquoi une maquette au dixième n’a qu’un millième
du volume initial. Les formules des autres solides usuels vérifient
le même facteur lorsque toutes leurs longueurs subissent le même changement.
\end{demonstration}
\editionRetenir{Le nombre de facteurs de longueur explique la puissance : 1 pour une longueur, 2 pour une aire, 3 pour un volume.}

\end{blocEdition}
